\documentclass[11pt,a4paper]{article}
\usepackage[utf8]{inputenc}
\usepackage[vietnamese]{babel}
\usepackage{amsmath,amssymb}
\usepackage{geometry}
\geometry{a4paper, total={180mm,260mm}, left=15mm, right=15mm, top=18mm, bottom=18mm}
\usepackage{tikz}
\usetikzlibrary{shapes.geometric, arrows.meta, positioning, calc}
\usepackage{tcolorbox}
\tcbuselibrary{skins,breakable}
\usepackage{xcolor}
\usepackage{multicol}
\usepackage{array}
\usepackage{fontawesome}
\usepackage{fancyhdr}

\pagestyle{fancy}
\fancyhf{}
\renewcommand{\headrulewidth}{0.5pt}
\renewcommand{\footrulewidth}{0.5pt}
\fancyhead[L]{\small\bfseries\color{teal}\faBook\ SỔ TAY KIẾN THỨC TOÁN 3}
\fancyhead[R]{\small\bfseries\color{teal}Lớp Toán Thầy Lĩnh -- loptoanthaylinh.vn}
\fancyfoot[C]{\small\color{gray}-- Sổ tay kiến thức: Bài 27. Giảm một số đi một số lần --}

\definecolor{primary}{HTML}{005F73}
\definecolor{secondary}{HTML}{0A9396}
\definecolor{accent}{HTML}{EE9B00}
\definecolor{danger}{HTML}{AE2012}
\definecolor{boxbg}{HTML}{F4F6F8}

\begin{document}

% --- BANNER TIÊU ĐỀ ---
\begin{tcolorbox}[
  enhanced,
  colback=white,
  colframe=primary,
  arc=3mm,
  boxrule=1.2pt,
  left=4mm, right=4mm, top=2.5mm, bottom=2.5mm,
  drop shadow=black!10!white
]
  \centering
  {\Large\bfseries\color{primary}\faBookmark\ BÀI 27: GIẢM MỘT SỐ ĐI MỘT SỐ LẦN}\\[1.5mm]
  {\small\itshape\color{black!75}Quy tắc giảm đi một số lần và phân biệt tuyệt đối với giảm đi một số đơn vị trên 1 trang A4}
\end{tcolorbox}

\vspace{-2mm}


\noindent
\begin{minipage}[t]{0.49\textwidth}\vspace{0pt}
\begin{tcolorbox}[
  enhanced,
  colback=blue!4!white,
  colframe=blue!75!black,
  arc=2mm, boxrule=0.8pt,
  left=2mm, right=2mm, top=1.8mm, bottom=1.8mm,
  title={\footnotesize\bfseries\color{white}\faKey\ 1. Quy tắc Giảm đi một số lần}
]
{\footnotesize
\begin{itemize}\setlength{\itemsep}{1.5pt}\setlength{\leftskip}{-2.5mm}
  \item Muốn \textbf{giảm một số đi nhiều lần}, ta lấy số đó \textbf{chia cho số lần}:
  $$\text{Số mới} = \text{Số ban đầu} : \text{Số lần}$$
  \item \textbf{Ví dụ cốt lõi:}
  \begin{itemize}\setlength{\itemsep}{1pt}\setlength{\leftskip}{-2.5mm}
    \item Giảm $18$ đi $3$ lần: $18 : 3 = 6$.
    \item Giảm $40$ đi $5$ lần: $40 : 5 = 8$.
    \item Giảm $72$ đi $8$ lần: $72 : 8 = 9$.
  \end{itemize}
  \item Giảm đi $2$ lần còn gọi là tìm \textbf{một nửa} của số đó (chia cho $2$).
\end{itemize}
}
\end{tcolorbox}
\end{minipage}
\hfill
\begin{minipage}[t]{0.49\textwidth}\vspace{0pt}
\begin{tcolorbox}[
  enhanced,
  colback=teal!4!white,
  colframe=teal!75!black,
  arc=2mm, boxrule=0.8pt,
  left=2mm, right=2mm, top=1.8mm, bottom=1.8mm,
  title={\footnotesize\bfseries\color{white}\faBalanceScale\ 2. So sánh hai mô hình giảm}
]
{\footnotesize
\begin{itemize}\setlength{\itemsep}{1.5pt}\setlength{\leftskip}{-2.5mm}
  \item \textbf{Mô hình 1: Giảm đi $n$ lần}
  $$\implies \text{Làm phép \textbf{CHIA}} \quad (a : n)$$
  \item \textbf{Mô hình 2: Giảm đi $n$ đơn vị}
  $$\implies \text{Làm phép \textbf{TRỪ}} \quad (a - n)$$
  \item \textbf{Bảng đối chiếu trực quan:}
  \begin{center}
  \renewcommand{\arraystretch}{1.15}
  \begin{tabular}{|c|c|c|}
    \hline
    \textbf{Số gốc} & \textbf{Giảm 4 lần} & \textbf{Giảm 4 đơn vị} \\
    \hline
    $24$ & $24 : 4 = \mathbf{6}$ & $24 - 4 = \mathbf{20}$ \\
    \hline
    $32$ & $32 : 4 = \mathbf{8}$ & $32 - 4 = \mathbf{28}$ \\
    \hline
  \end{tabular}
  \end{center}
\end{itemize}
}
\end{tcolorbox}
\end{minipage}

\vspace{-2.5mm}
% --- KHỐI 3: SƠ ĐỒ TƯ DUY TIKZ ---
\noindent
\begin{tcolorbox}[
  enhanced,
  colback=orange!4!white,
  colframe=orange!80!black,
  arc=2mm, boxrule=0.8pt,
  left=2mm, right=2mm, top=1.5mm, bottom=1.5mm,
  title={\footnotesize\bfseries\color{white}\faSitemap\ SƠ ĐỒ TƯ DUY: BẢN ĐỒ TÁC TOÁN TĂNG / GIẢM TRONG TOÁN 3}
]
\centering
\begin{tikzpicture}[
  box/.style={draw, rounded corners=2pt, align=center, font=\scriptsize\bfseries, inner sep=2.5pt, minimum height=0.7cm},
  arr/.style={-{Stealth[length=2mm]}, thick}
]
  \node[box, fill=purple!15, draw=purple!80!black, text width=2.6cm] (c1) at (0, 0) {SỐ BAN ĐẦU\\Đại lượng $A$};
  \node[box, fill=teal!15, draw=teal!80!black, text width=2.8cm] (c2) at (4.0, 0.7) {GẤP $n$ LẦN\\$\mathbf{A \times n}$\\(Phép nhân)};
  \node[box, fill=blue!12, draw=blue!80!black, text width=2.8cm] (c3) at (4.0, -0.7) {GIẢM $n$ LẦN\\$\mathbf{A : n}$\\(Phép chia)};
  \node[box, fill=green!15, draw=green!70!black, text width=2.8cm] (c4) at (8.5, 0.7) {TĂNG $n$ ĐƠN VỊ\\$\mathbf{A + n}$\\(Phép cộng)};
  \node[box, fill=red!12, draw=red!70!black, text width=2.8cm] (c5) at (8.5, -0.7) {GIẢM $n$ ĐƠN VỊ\\$\mathbf{A - n}$\\(Phép trừ)};

  \draw[arr, teal!80!black] (c1) -- (c2);
  \draw[arr, blue!80!black] (c1) -- (c3);
  \draw[arr, green!70!black] (c2) -- (c4);
  \draw[arr, red!70!black] (c3) -- (c5);
\end{tikzpicture}
\end{tcolorbox}

\vspace{-2.5mm}
% --- KHỐI 4: VÍ DỤ & MẸO TRÁNH BẪY ---
\noindent
\begin{minipage}[t]{0.49\textwidth}\vspace{0pt}
\begin{tcolorbox}[
  enhanced,
  colback=green!4!white,
  colframe=green!60!black,
  arc=2mm, boxrule=0.8pt,
  left=2mm, right=2mm, top=1.8mm, bottom=1.8mm,
  title={\footnotesize\bfseries\color{white}\faCheckCircle\ Ví dụ thực hành chuẩn mực}
]
{\footnotesize
\textbf{1. Giải toán có lời văn:}\\
Đoạn thẳng $AB$ dài $24\text{ cm}$. Đoạn thẳng $CD$ ngắn hơn đoạn $AB$ giảm đi $3$ lần. Độ dài đoạn thẳng $CD$ là:\\
\phantom{Lời giải: }$24 : 3 = \mathbf{8}\text{ (cm)}$.\\
\textbf{2. Bài toán kết hợp:}\\
Mẹ có $30$ quả cam, mẹ đem biếu bà giảm đi $5$ lần số cam ban đầu. Hỏi mẹ còn lại bao nhiêu quả cam?\\
Số cam biếu bà: $30 : 5 = 6$ (quả).\\
Số cam còn lại: $30 - 6 = \mathbf{24}$ (quả).
}
\end{tcolorbox}
\end{minipage}
\hfill
\begin{minipage}[t]{0.49\textwidth}\vspace{0pt}
\begin{tcolorbox}[
  enhanced,
  colback=red!4!white,
  colframe=red!70!black,
  arc=2mm, boxrule=0.8pt,
  left=2mm, right=2mm, top=1.8mm, bottom=1.8mm,
  title={\footnotesize\bfseries\color{white}\faExclamationTriangle\ Mẹo tránh bẫy trắc nghiệm}
]
{\footnotesize
\begin{itemize}\setlength{\itemsep}{1.5pt}\setlength{\leftskip}{-2.5mm}
  \item \textbf{Bẫy từ ngữ "lần" và "đơn vị":}\\
  Đọc thật chậm từ khóa trong câu hỏi:\\
  -- Thấy chữ \textbf{"giảm ... lần"} $\to$ làm tính \textbf{chia} ($:$).\\
  -- Thấy chữ \textbf{"giảm ... cm/kg/quả"} hoặc \textbf{"bớt đi"} $\to$ làm tính \textbf{trừ} ($-$).
  \item Rất nhiều bạn học sinh tính nhầm "giảm $15$ đi $3$ lần" thành $15 - 3 = 12$ thay vì $15 : 3 = 5$.
\end{itemize}
}
\end{tcolorbox}
\end{minipage}


\end{document}
